The J-track papers J.1--J.6 develop each method and the paradox theory
in depth.
We summarise the key finding from each paper.

\paragraph{J.1 (Huntington-Hill).}
Huntington-Hill's priority formula $P_i(s) = p_i/\sqrt{s(s+1)}$
implements the geometric mean rounding rule.
The paper proves that this minimises the relative difference in
per-capita representation across all pairs of states,
documents the 1941 statutory adoption (2~U.S.C.\ \S2a),
and analyses the Supreme Court's upholding of this statutory
choice in \textit{United States Dept.\ of Commerce v.\ Montana},
503~U.S.\ 442 (1992).
The SHA-verified 2020 apportionment appears in the appendix.

\paragraph{J.2 (Webster).}
Webster's arithmetic mean rounding is slightly more large-state-favoring
than Huntington-Hill's geometric mean.
For the 2020 Census the methods agree for 49 states;
the paper identifies the one state near the rounding boundary
and documents Webster's historical use in 1842 and 1911--1940.
Webster's optimality in terms of minimising squared deviation from quota
makes it the natural academic benchmark for evaluating HH's small-state
bias.

\paragraph{J.3 (Adams).}
Adams ceiling rounding maximises the minimum per-capita representation
across all states --- the most small-state-favorable divisor method.
Under 2020 data, Adams gives several small states an extra seat
at the expense of large states.
The paper proves paradox immunity for ceiling-rounding divisor methods
using the same monotonicity argument that applies to all divisor methods.

\paragraph{J.4 (Jefferson/D'Hondt).}
Jefferson's floor rounding is the most large-state-favoring divisor method
and, internationally, the most widely used electoral apportionment method.
The paper proves the formal equivalence between Jefferson and d'Hondt
(the priority formula $p_i/s$ applies to both state populations and
party vote totals), and documents Jefferson's use in the United States
from 1790 to 1840 and its eventual abandonment due to large-state bias.

\paragraph{J.5 (Paradoxes).}
The central paper of the J-track.
It proves susceptibility for Hamilton (with historical examples:
Alabama 1880, Population 1900, Oklahoma 1907) and immunity for all
four divisor methods.
The proof of divisor method immunity uses the monotone priority-queue
argument.
The Balinski-Young impossibility theorem \citep{balinski1982fair}
--- no method can simultaneously satisfy the quota rule and be
paradox-immune --- provides the theoretical capstone.

\paragraph{J.6 (Implementation).}
Documents the \textsc{bisect-apportion} crate in detail:
integer arithmetic for exact priority comparison, the SHA-256
verification protocol, the 61-test suite at L0/L1/L2 levels,
and comparative results for all five methods on 2020 Census data.
The crate's unified interface enables the comparative analysis
in J.0--J.5 to be reproduced computationally.
